\documentclass[11pt]{article}

\usepackage{amsmath,amssymb}
\usepackage{xurl}
\usepackage[colorlinks=true,linkcolor=blue,citecolor=blue,urlcolor=blue]{hyperref}

\input{command}

\title{Implicit Nilpotent-Space Classification in the Low-Dimensional Quantum Wielandt Corollary}
\author{}
\date{2026-08-24}

\begin{document}
\maketitle
\gapnote{open-gap}{wip}

\section{The assertion}

The corollary after Wolf's quantum Wielandt inequality states that a primitive
quantum channel on $\MN{d}$, for $d=2$ or $d=3$, has an element of its one-step
Kraus span with a nonzero eigenvalue and consequently has Wielandt index at
most $d^2$ \cite[Chapter~6, lines 1110--1126]{Wolf2012Quantum}.  Its proof
invokes a classification of nilpotent matrix spaces without proving the
classification or resolving the bibliography placeholders in the local source.
The same low-dimensional observation appears after Theorem~1 of
\cite{Sanz2010Quantum}.

For $d=2$, a nilpotent linear space has dimension at most one.  For $d=3$,
Wolf describes two similarity types of nilpotent spaces of dimension greater
than one: the strictly upper-triangular space and an exceptional
two-dimensional space.  The former is incompatible with trace preservation;
the latter gives a channel whose peripheral spectral radius is degenerate.
These classification and obstruction arguments are not contained in the
lecture notes.  Fasoli's classification work is a relevant external source
\cite{Fasoli1997Classification}, but the exact chain needed for Wolf's stated
corollary must still be reconstructed and checked.

\section{Formalized part}

The declaration
\leanid{QICLean.finrank\_le\_one\_of\_isNilpotent\_submodule\_fin\_two}
proves the complete $d=2$ nilpotent-space dimension bound, and
\leanid{QICLean.exists\_common\_ker\_of\_isNilpotent\_submodule\_fin\_two}
proves the stronger common-kernel form used for channels.  Combining this with
trace preservation,
\leanid{Kraus.exists\_nonzero\_eigenvector\_mem\_wordSpan\_one\_fin\_two}
proves that the one-step span of every trace-preserving Kraus family on
$\MN{2}$ contains an element with a nonzero eigenvalue.

For $d=3$,
\leanid{QICLean.exists\_common\_ker\_of\_forall\_sq\_eq\_zero}
proves the square-zero branch, and
\leanid{Kraus.exists\_mem\_wordSpan\_one\_sq\_ne\_zero\_fin\_three}
excludes that branch for a trace-preserving family.  The remaining
nilpotency-index-three branch, including the exceptional similarity type and
its primitivity obstruction, is not yet formalized.

\section{Remaining gap}

A source-faithful proof still requires two ingredients.

\begin{enumerate}
\item Complete the $3\times3$ classification by proving that a nilpotent space
containing an element of nilpotency index three either has a common kernel or
is similar to the exceptional two-dimensional space, and prove that this
exceptional Kraus pattern is incompatible with Wolf's primitivity condition.
\item Supply in QICLean the conditional implication from a nonzero eigenvalue
in the one-step span to full word span at length $d^2$.  The current QICLean
checkout contains the eigenvector-spreading component but no assembled bound;
the numbering-coverage document currently records Wolf's Theorem~6.9 as
living in the downstream tensor-network layer.
\end{enumerate}

Until both ingredients are present, the source corollary must not carry a
\texttt{\textbackslash leanok} marker in the QICLean blueprint.  This work is
tracked by \ghissue{65}.

\bibliographystyle{plain}
\bibliography{references}

\end{document}
